\documentclass{article}
\usepackage[T1]{fontenc}
\usepackage{lmodern}
\usepackage{graphicx}
\usepackage{amsmath,amssymb}
\usepackage[a4paper,margin=2cm]{geometry}
\usepackage[absolute,overlay]{textpos}
\setlength{\TPHorizModule}{1mm}
\setlength{\TPVertModule}{1mm}
\usepackage[indonesian]{babel}
\usepackage{hyperref}
\setlength{\emergencystretch}{2em}
% unit: Halaman 44

\begin{document}

% === Halaman 44 (python/native) ===

44

\# Program utama

\#Alice membangkitkan kunci publik dan kunci privatnya

print("\textbackslash{}nKUNCI PUBLIK DAN KUNCI PRIVAT ALICE:")

PembangkitanKunciRSA('kunci\_publik\_Alice.pem', 'kunci\_privat\_Alice.pem')

\#Bob membangkitkan kunci publik dan kunci privatnya

print("\textbackslash{}nKUNCI PUBLIK DAN KUNCI PRIVAT BOB:")

PembangkitanKunciRSA("kunci\_publik\_Bob.pem", "kunci\_privat\_Bob.pem")

\#Bob mengirim pesan kepada Alice

pesan = input("\textbackslash{}nBob, ketikkan pesan teks buat Alice: ")

\#Bob mengenkripsi pesan dengan menggunakan kunci publik Alice

ciphertext = encrypt\_message(pesan, 'kunci\_publik\_Alice.pem')

print("\textbackslash{}nCiphertext (dari Bob):", ciphertext.hex())

\#Alice mendekripsi ciphertext dari Bob menggunakab kunci privatnya sendiri 

(kunci privat Alice)

decrypted\_message = decrypt\_message(ciphertext, 'kunci\_privat\_Alice.pem')

print("\textbackslash{}nPesan setelah dekripsi oleh Alice:", decrypted\_message)

\end{document}
